\chapter{Hecke, Adams, and Witt operations}
\label{v4-arith:hecke:chapter}

An operation on mode indices does not determine its arithmetic meaning.
Divisor substitutions, index shifts and Hecke operators act on
different structures, even when their first values agree.
Their composition laws separate them. Integral lattices, distinguished
states and completed domains give further distinctions.

We begin with three explicit formulas. The divisor substitutions have
a noncommutative algebra of boundary corrections. A second divisor
formula gives a genuine linear Hecke action. The index shifts preserve
the integral symmetric-function lattice and satisfy additional Frobenius congruences.
Big Witt coordinates retain this integral information when ghost values
lose it. Finally, the radial transform and weighted operator domains
specify which comparisons survive Hilbert completion.

\section{Three operations on mode indices}
\label{v4-arith:hecke:sec:three}

Let $R$ be a commutative unital ring and $S$ a finite set of species.
Define the free module and its symmetric algebra by
\[
 M_R=\bigoplus_{s\in S,\ k\geq1}R x_{s,k},\qquad
 P_R=\operatorname{Sym}_R(M_R).
\]
These objects have cohomological degree zero. Each generator has
particle degree one and weight $k$.
The unit $1\in P_R$ is the
polynomial vacuum. All polynomials have finite support.

The species index and the index of an operation are independent.

For $n\geq1$, define two unital $R$-algebra endomorphisms by
\begin{equation}
 \label{v4-arith:hecke:eq:substitutions}
 D_nx_{s,k}=\sum_{d\mid(n,k)}d\,x_{s,nk/d},\qquad
 \psi_nx_{s,k}=x_{s,nk}.
\end{equation}
The first is the \emph{divisor substitution}. The second is the
\emph{index shift}. The sums are finite and use integer coefficients.
Thus both definitions commute with scalar extension.

For an arbitrary species set, take the union over its finite subsets.
For $S=\varnothing$, both maps are the identity on $R$.

A third formula defines an $R$-linear endomorphism of $M_R$:
\begin{equation}
 \label{v4-arith:hecke:eq:linear}
 \mathsf H_nx_{s,k}=\sum_{d\mid(n,k)}d\,x_{s,nk/d^2}.
\end{equation}
The index is integral because $d$ divides both $n$ and $k$.
The squared denominator changes the composition law.
We suppress the species when examining one summand.

\begin{example}
\label{v4-arith:hecke:ex:first-mode}
For a prime $p$, all three maps send $x_1$ to $x_p$.
Their next values are
\[
 D_px_p=x_{p^2}+px_p,\qquad
 \psi_px_p=x_{p^2},\qquad
 \mathsf H_px_p=x_{p^2}+px_1.
\]
Consequently
\[
 (D_p^2-D_{p^2})x_1=px_p,\qquad
 (D_p^2-D_{p^2}-pI)x_1=p(x_p-x_1).
\]
Both differences are nonzero over $\mathbb Z$.
The index shifts satisfy $\psi_m\psi_n=\psi_{mn}$ on every generator,
hence on $P_R$. The linear Hecke relation will instead include
the additional term $pI$.
\end{example}

Extending a linear generator map $A:M_R\to M_R$ to
$\operatorname{Sym}(A):P_R\to P_R$ preserves composition.
It does not preserve addition as a function of $A$.
Already
\[
 \operatorname{Sym}(2A)(x^2)=4(Ax)^2,
 \qquad 2\operatorname{Sym}(A)(x^2)=2(Ax)^2.
\]
Moreover, every such extension fixes the vacuum. Over $\mathbb Z$,
the putative polynomial Hecke relation
\[
 \operatorname{Sym}(\mathsf H_p)^2
 =\operatorname{Sym}(\mathsf H_{p^2})+pI
\]
already fails on $1$, where its two sides give $1$ and $1+p$.
Linear operator relations and multiplicative substitutions must
therefore retain their different carriers.

\section{Composition and the boundary of a prime chain}
\label{v4-arith:hecke:sec:composition}

Fix a prime $p$ and $u\geq1$ with $(u,p)=1$.
Write $e_r=x_{up^r}$ for $r\geq0$, and let $A_a$ be the
restriction of $D_{p^a}$ to this free submodule.
Then $A_0=I$ and
\begin{equation}
 \label{v4-arith:hecke:eq:prime-chain}
 A_ae_r=\sum_{i=0}^{\min(a,r)}p^i e_{r+a-i}.
\end{equation}
The lower endpoint $r=0$ truncates the sum.
Farther along the chain, the same operators act by polynomials
in the forward shift. Their noncommutativity is confined to
the truncated region.

\begin{theorem}[Prime-chain composition]
\label{v4-arith:hecke:thm:composition}
For $a,b,r\geq0$, put $[v]_+=\max(v,0)$ and
\[
 N_{a,b}(r,t)=\mathbf1_{t\leq r+b}
 \bigl[\min(b,r,t)-\max(0,t-a)+1\bigr]_+.
\]
Over every commutative coefficient ring,
\begin{equation}
 \label{v4-arith:hecke:eq:composition}
 A_aA_be_r=\sum_{t=0}^{a+b}p^tN_{a,b}(r,t)e_{r+a+b-t}.
\end{equation}
For $1\leq a<b$, the boundary commutator is
\begin{equation}
 \label{v4-arith:hecke:eq:commutator}
 [A_a,A_b]e_1=-p^{a+1}e_b.
\end{equation}
Thus, over a nonzero characteristic-zero ring, $A_a$ and $A_b$
commute exactly when $a=b$ or $ab=0$.
\end{theorem}
\begin{proof}
A term of $A_aA_be_r$ chooses integers
\[
 0\leq i\leq\min(b,r),\qquad
 0\leq j\leq\min(a,r+b-i).
\]
Its coefficient is $p^{i+j}$ and its output index is $r+a+b-i-j$.
Fixing $t=i+j$ leaves the conditions
\[
 t\leq r+b,\qquad \max(0,t-a)\leq i\leq\min(b,r,t).
\]
Their number is $N_{a,b}(r,t)$, proving the first formula.

For the commutator, first use $A_be_1=e_{b+1}+pe_b$.
Applying $A_a$ gives successive coefficients
$1,2p,\ldots,2p^a,p^{a+1}$.
In the reverse order, $A_ae_1=e_{a+1}+pe_a$.
Since $b\geq a+1$, its final coefficient is $2p^{a+1}$.
All other coefficients agree. Their difference proves
\eqref{v4-arith:hecke:eq:commutator}.

The nonzero integer $p^{a+1}$ remains nonzero in a characteristic-zero
ring. The commuting cases are immediate, and exchanging $a,b$
covers the other order.
\end{proof}

For example, the entire first commutator is
\[
 [A_1,A_2]e_r=
 \begin{cases}-p^2e_2,&r=1,\\0,&r\ne1.\end{cases}
\]
For $r\geq2$, both operators have their untruncated polynomial form.
The two remaining basis vectors give the displayed boundary term.

\begin{corollary}[Independent prime coordinates]
\label{v4-arith:hecke:cor:global}
Set $a_p=v_p(m)$, $b_p=v_p(n)$ and $r_p=v_p(k)$.
Then
\[
 D_mD_nx_k=
 \sum_{(t_p)}\left(\prod_{p\mid mn}p^{t_p}
 N_{a_p,b_p}(r_p,t_p)\right)
 x_{mnk/\prod_{p\mid mn}p^{t_p}},
 \qquad 0\leq t_p\leq a_p+b_p.
\]
Terms with zero coefficient are omitted.
In particular, $D_mD_n=D_{mn}$ when $(m,n)=1$.
For a nonempty species set over a nonzero characteristic-zero ring,
\begin{align*}
 D_mD_n=D_{mn}&\quad\Longleftrightarrow\quad(m,n)=1,\\
 [D_m,D_n]=0&\quad\Longleftrightarrow\quad
 \text{for every }p,\ a_pb_p=0\text{ or }a_p=b_p.
\end{align*}
Over $\mathbb F_\ell$, omit the condition at $p=\ell$ in both criteria.
\end{corollary}
\begin{proof}
Multiplication or division by a power of $p$ leaves every other
prime valuation unchanged. The divisor choices therefore factor
over primes, giving the formula and the coprime statement.
Equality on generators extends multiplicatively to $P_R$.

If $a_p,b_p>0$, then
\[
 A_{a_p}A_{b_p}e_0
 =\sum_{j=0}^{\min(a_p,b_p)}p^je_{a_p+b_p-j}.
\]
Its $j=1$ term is absent from $A_{a_p+b_p}e_0$.
Selecting the highest terms at the other primes isolates its
nonzero coefficient. This proves the first necessity.

For the second, use $x_p$ at a prime with unequal positive exponents.
At other primes, both orders agree on their initial vector.
Their highest terms and \eqref{v4-arith:hecke:eq:commutator}
again isolate a nonzero power of $p$.

In characteristic $\ell$, the terms with positive powers of $\ell$
vanish, so $D_{\ell^a}=\psi_{\ell^a}$ on $P_{\mathbb F_\ell}$.
At every other prime, the same powers of $p$ remain nonzero.
This proves precisely the stated exception.
\end{proof}

The index shifts interact with the divisor substitutions by
\[
 \psi_mD_nx_k=\sum_{d\mid(n,k)}d\,x_{mnk/d},\qquad
 D_n\psi_mx_k=\sum_{d\mid(n,mk)}d\,x_{mnk/d}.
\]
In characteristic zero they commute exactly when $(m,n)=1$.
Coprimality equates the divisor conditions. If a common prime exists,
$k=1$ supplies an extra term with a nonzero integer coefficient.

There is also a monoid action retaining the prime divisor maps.
Define $\Theta_n=\prod_pD_p^{v_p(n)}$.
Different prime maps commute, so $\Theta_m\Theta_n=\Theta_{mn}$.

For $a\geq1$, its prime-chain values are
\[
 A_1^ae_r=\sum_{j=0}^a\binom ajp^je_{r+a-j}\quad(r\geq1),
 \qquad
 A_1^ae_0=\sum_{j=0}^{a-1}\binom{a-1}jp^je_{a-j}.
\]
The first formula is the binomial theorem on the invariant tail.
The second follows after $e_0\mapsto e_1$.
This action changes the higher operations. Its prime maps retain
the Frobenius obstruction computed below.

\section{The algebra of linear divisor operators}
\label{v4-arith:hecke:sec:presentation}

The prime-chain algebra admits a finite test for each proposed relation.
Each test controls a finite boundary region and the entire infinite tail.

Let $Ue_r=e_{r+1}$, $Be_0=0$ and $Be_r=e_{r-1}$ for $r\geq1$.
Then $BU=I$ and $U^iB^i$ projects onto the tail $r\geq i$.
Consequently
\[
 A_a=U^a\sum_{i=0}^ap^iB^i,
 \qquad A_{a+1}=UA_a+p^{a+1}U^{a+1}B^{a+1}.
\]
These identities place the generated algebra inside the algebra of
the shift and its partial inverse. They need not generate that
entire ambient algebra.

Put $q_a(z)=\sum_{i=0}^ap^iz^{a-i}$ for $a\geq1$.
For a word $w=Z_{a_l}\cdots Z_{a_1}$, define
\[
 d(w)=\sum_ja_j,\qquad q_w=\prod_jq_{a_j},
\]
and
\[
 B_w(r,s)=
 \sum_{\substack{0\leq i_j\leq a_j,\ \sum_j(a_j-i_j)=s\\
 i_j\leq r+\sum_{h<j}(a_h-i_h)\ \text{for each }j}}
 p^{\sum_ji_j}.
\]
The rightmost factor acts first. For the empty word, set $d(w)=0$,
$q_w=1$ and $B_w(r,s)=\mathbf1_{s=0}$.

\begin{theorem}[Tail and boundary presentation]
\label{v4-arith:hecke:thm:presentation}
Let $F=\sum_wc_ww$ be a finite element of
$R\langle Z_1,Z_2,\ldots\rangle$, and set $d=\max_wd(w)$.
For $F=0$, take $d=0$.
Then $F(A_1,A_2,\ldots)=0$ if and only if
\begin{align*}
 &\sum_wc_wq_w(z)=0\quad\text{in }R[z],\\
 &\sum_wc_wB_w(r,s)=0
       \quad(0\leq r<d,\ 0\leq s\leq d).
\end{align*}
These conditions define the kernel ideal $I_{p,R}$, giving
\[
 R[A_1,A_2,\ldots]\simeq
 R\langle Z_1,Z_2,\ldots\rangle/I_{p,R}.
\]
Every finite expression has a finite coefficient test.
No finite generating set for the ideal is asserted.
\end{theorem}
\begin{proof}
Expanding the successive divisor choices gives
$w(A)e_r=\sum_sB_w(r,s)e_{r+s}$.
For $r\geq d$, every intermediate index is at least $r$.
All choices $0\leq i_j\leq a_j$ are then allowed.
Their coefficients are exactly those of $q_w(z)$.
The polynomial identity therefore tests all vectors with $r\geq d$.

The remaining equations test every earlier basis vector.
Together they are precisely the kernel of evaluation $Z_a\mapsto A_a$.
The kernel is a two-sided ideal, and the quotient is its image.
\end{proof}

Equivalently, every evaluated expression is uniquely $q(U)+E$.
Here $E$ is a finite matrix, supported in columns $r<d$ and rows
$r\leq r+s\leq r+d$. Uniqueness follows by testing arbitrarily
large $e_r$. The tail-symbol map onto $R[z]$ is surjective,
since $q_1=z+p$. Its kernel consists of the finite matrices
which occur in the generated algebra.
The commutator $[A_1,A_2]$ gives a nonzero example over $\mathbb Z$.

For finitely many primes, use a basis indexed by their valuation vectors.
Word coefficients are products of the corresponding prime-chain coefficients.
In each coordinate, separate the finitely many boundary values from
the unrestricted tail. This gives finitely many mixed boundary tests
and polynomial identities for every finite expression.
Every expression in the full prime family involves only finitely many primes.

\section{The linear Hecke law and its determinant}
\label{v4-arith:hecke:sec:linear}

The partial inverse enters differently in
\eqref{v4-arith:hecke:eq:linear}. On $M_R$, define
\[
 U_px_k=x_{pk},\qquad
 B_px_k=\begin{cases}x_{k/p},&p\mid k,\\0,&p\nmid k.\end{cases}
\]
Then $B_pU_p=I$, while $U_pB_p$ projects onto the
$p$-divisible indices. At distinct primes all four pairs commute.
The linear operators satisfy
\[
 \mathsf H_{p^r}=\sum_{j=0}^rp^jU_p^{r-j}B_p^j.
\]

Multiplying by $U_p+pB_p$ gives
\[
 \mathsf H_p\mathsf H_{p^r}
 =\mathsf H_{p^{r+1}}+p\mathsf H_{p^{r-1}}
 \qquad(r\geq1).
\]
Indeed, the final backward-shift term supplies $p^{r+1}B_p^{r+1}$.
Every other backward-shift term cancels a forward shift using
$B_pU_p=I$, leaving $p\mathsf H_{p^{r-1}}$.

\begin{proposition}[Universal determinant identity]
\label{v4-arith:hecke:prop:determinant}
For independent indeterminates $T,d$, define
\[
 H_0=1,\qquad H_1=T,\qquad H_{r+1}=TH_r-dH_{r-1}\quad(r\geq1).
\]
Then in $\mathbb Z[T,d]$,
\[
 H_rH_s=\sum_{j=0}^{\min(r,s)}d^jH_{r+s-2j}.
\]
The identity also holds after substituting any commuting pair of
endomorphisms of a module over a commutative ring.
\end{proposition}
\begin{proof}
Substitute $T=u+v$ and $d=uv$.
The recurrence gives $H_r=\sum_{i=0}^ru^iv^{r-i}$.
For $r\leq s$, the coefficient of $u^iv^{r+s-i}$ in the product is
\[
 1+\min(r,i,r+s-i)\qquad(0\leq i\leq r+s).
\]
The right side has one contribution for each
$0\leq j\leq\min(r,i,r+s-i)$, so the coefficients agree.

The substitution is injective: in lexicographic order with $u>v$,
$(u+v)^a(uv)^b$ has leading monomial $u^{a+b}v^b$.
These leading monomials are distinct for distinct pairs $(a,b)$.
The identity therefore holds before specialization, including $r=s$.
Evaluation at commuting endomorphisms preserves it.
\end{proof}

\begin{corollary}
\label{v4-arith:hecke:cor:linear-law}
Over every commutative coefficient ring,
\[
 \mathsf H_m\mathsf H_n
   =\sum_{d\mid(m,n)}d\,\mathsf H_{mn/d^2}.
\]
\end{corollary}
\begin{proof}
At each prime, specialize the proposition to $(\mathsf H_p,pI)$.
The divisor formula factors $\mathsf H_n$ over its prime powers.
Multiplication of the local identities proves the global formula.
\end{proof}

This is a representation of the stated linear operator algebra.
The algebra generated by all $U_p,B_p$ also retains these relations
in its regular representation.
Neither construction identifies its multiplication with the vertex
products of a Heisenberg algebra.

\section{Satake parameters and the radial state}
\label{v4-arith:hecke:sec:satake}

The determinant in the universal identity has a geometric normalization.
Fix $K_p=\mathrm{GL}_2(\mathbb Z_p)$ and give it Haar volume one.
The complex spherical Hecke algebra is
\[
 \mathcal H_p=C_c(K_p\backslash\mathrm{GL}_2(\mathbb Q_p)/K_p,\mathbb C),
\]
The product is convolution. The two characteristic functions used below are
\[
 T_p=\mathbf1_{K_p\operatorname{diag}(p,1)K_p},\qquad
 \mathsf z_p=\mathbf1_{K_p\operatorname{diag}(p,p)K_p}.
\]

\begin{theorem}[Normalized spherical Satake isomorphism]
\label{v4-arith:hecke:thm:satake}
The normalized Satake isomorphism identifies
\[
 \mathcal H_p\simeq
 \mathbb C[X^{\pm1},Y^{\pm1}]^{S_2},\qquad
 T_p\longmapsto\sqrt p(X+Y),\qquad \mathsf z_p\longmapsto XY.
\]
At a character, put $\alpha=\sqrt pX$ and $\beta=\sqrt pY$.
Then $\alpha\beta=p\mathsf z_p$, using the scalar central character on the right.
Trivial central character gives $\mathsf z_p=1$.
Temperedness is the additional condition $|X|=|Y|=1$.
\end{theorem}

This is the convention of Gross, \emph{On the Satake isomorphism},
Proposition~3.6, equations~(3.13)--(3.14), and Section~8.
In particular, trivial central character does not imply temperedness.
The higher elements $H_{p^r}\in\mathcal H_p$ are specified by
the complete symmetric Satake polynomials:
\[
 \sum_{r\geq0}H_{p^r}t^r=(1-T_pt+p\mathsf z_pt^2)^{-1}.
\]
This follows by multiplying $(1-\sqrt pXt)$ and $(1-\sqrt pYt)$
in the Satake algebra and using injectivity.
These higher elements need not be individual double-coset characteristic functions.

For a fixed pair $\alpha\beta=p$, the diagonal matrix
$S_p=\operatorname{diag}(\alpha,\beta)$ gives
\[
 \operatorname{Tr}(\operatorname{Sym}^rS_p)
 =\sum_{j=0}^r\alpha^{r-j}\beta^j=H_r(\alpha+\beta,p).
\]
This follows from the monomial eigenbasis of the symmetric power.
The matrix, its scalar traces and the operators $\mathsf H_n$
are three separate representations of the parameter formulas.
In a basis $a,b$, the line $\mathbb C(a+b)$ is $S_p$-invariant
exactly when $\alpha=\beta$. Its image is
$\alpha a+\beta b$, which proves the criterion.

If the pair is tempered, write
\[
 \{\alpha,\beta\}
 =\{\sqrt p\,e^{i\theta_p},\sqrt p\,e^{-i\theta_p}\}.
\]
The angle is defined modulo $2\pi$ and sign.

For any family of such angles, the product
\[
 L_\theta(s)=\prod_p
 \frac1{1-2\cos\theta_p\,p^{-s}+p^{-2s}}
\]
converges absolutely and locally uniformly for $\operatorname{Re}s>1$.
Indeed, its logarithm is the absolutely convergent sum
\[
 \sum_p\sum_{m\geq1}\frac{2\cos(m\theta_p)}m p^{-ms}.
\]
On every closed smaller half-plane, its absolute values are bounded
by a constant times $\sum_pp^{-\operatorname{Re}s}$.

This constructs the local-parameter Euler product on that domain.
An automorphic identification or continuation beyond it requires
additional arithmetic data. Replacing unitary parameters by
$\alpha,\beta$ shifts the Dirichlet variable by one half.

\subsection{The integral cyclic module of a tree}

Fix a rooted $(p+1)$-regular tree $(\mathcal T_p,o)$.
Use counting measure on its vertices and let $S_r$ be the sphere
of radius $r$ about $o$. Its size is
\[
 s_0=1,\qquad s_r=(p+1)p^{r-1}\quad(r\geq1).
\]
The radial Hilbert space $\mathcal R_p$ consists of square-summable
functions constant on each sphere. Its orthonormal basis is
$e_r=s_r^{-1/2}\mathbf1_{S_r}$.

Adjacency is $(Tf)(v)=\sum_{w\sim v}f(w)$.
Cauchy--Schwarz bounds its norm on the vertex Hilbert space by $p+1$,
and symmetry makes it self-adjoint. Counting edges gives
\begin{align*}
 Te_0&=\sqrt{p+1}\,e_1,\\
 Te_1&=\sqrt{p+1}\,e_0+\sqrt p\,e_2,\\
 Te_r&=\sqrt p(e_{r-1}+e_{r+1})\quad(r\geq2).
\end{align*}
Distance defines a filtration, not a grading preserved by $T$.

Define monic integer polynomials by
\[
 R_0=1,\quad R_1=x,\quad R_2=x^2-(p+1),\qquad
 R_{r+1}=xR_r-pR_{r-1}\quad(r\geq2).
\]
The recurrence gives $R_r(T)e_0=\mathbf1_{S_r}$.
Monicity then proves an integral cyclic-module isomorphism
\begin{equation}
 \label{v4-arith:hecke:eq:integral-tree}
 \mathbb Z[x]\xrightarrow{\sim}
 \bigoplus_{r\geq0}\mathbb Z\mathbf1_{S_r},\qquad
 f\longmapsto f(T)e_0.
\end{equation}
No prime is inverted in this statement.
The normalized Hilbert basis instead uses real square roots.

\begin{theorem}[Radial transform with its state]
\label{v4-arith:hecke:thm:radial}
Put $x(\theta)=2\sqrt p\cos\theta$ for $0\leq\theta\leq\pi$, and define
\[
 d\nu_p(\theta)=
 \frac{2p(p+1)\sin^2\theta}
 {\pi(p^2+1-2p\cos2\theta)}\,d\theta.
\]
The measure has mass one. The assignment
\[
 \mathcal U_pe_r=\frac{R_r(x(\theta))}{\sqrt{s_r}}
\]
extends to a unitary $\mathcal R_p\xrightarrow{\sim}L^2([0,\pi],\nu_p)$.
It sends $e_0$ to $1$ and intertwines $T$ with multiplication by $x$.
In particular, $\|T|_{\mathcal R_p}\|=2\sqrt p$, and
\begin{equation}
 \label{v4-arith:hecke:eq:radial-state}
 \langle e_0,f(T)e_0\rangle
 =\int_0^\pi f(2\sqrt p\cos\theta)\,d\nu_p(\theta)
\end{equation}
for every bounded Borel function on the spectral interval.
\end{theorem}
\begin{proof}
The tridiagonal recurrence makes $e_0$ cyclic.
After deleting its row and column, the remaining Jacobi operator
has every off-diagonal coefficient equal to $\sqrt p$.
Its first diagonal resolvent entry satisfies
\[
 g(z)=\frac1{z-pg(z)}
 =\frac{z-\sqrt{z^2-4p}}{2p}.
\]
The branch is analytic off $[-2\sqrt p,2\sqrt p]$ and asymptotic to $z$.
It is selected by $g(z)\sim z^{-1}$.

Taking the Schur complement at the original root gives
\[
 m(z)=\langle e_0,(z-T)^{-1}e_0\rangle
 =\frac1{z-(p+1)g(z)}.
\]
For $|x|<2\sqrt p$, Stieltjes inversion yields
\[
 -\frac1\pi\operatorname{Im}m(x+i0)
 =\frac{(p+1)\sqrt{4p-x^2}}{2\pi((p+1)^2-x^2)}.
\]

A pole outside the interval would require $g^2=1$ and $z=\pm(p+1)$.
The chosen branch instead has $g(\pm(p+1))=\pm p^{-1}$.
At either endpoint the denominator is nonzero because $p>1$.
Thus $m$ has continuous finite boundary values on the interval
and analytic continuation across its complement.

The inversion formula gives only the displayed absolutely continuous measure.
Its mass is one by $m(z)\sim z^{-1}$.
Changing variables to $x=2\sqrt p\cos\theta$ gives $\nu_p$.

The spectral theorem and \eqref{v4-arith:hecke:eq:integral-tree} imply
\[
 \int R_r(x)R_s(x)\,d\nu_p
 =\langle\mathbf1_{S_r},\mathbf1_{S_s}\rangle=s_r\delta_{rs}.
\]
The normalized polynomials are therefore orthonormal.
They span all polynomials by monicity.
Polynomial approximation on the compact interval proves their
density in $L^2$, hence the asserted unitary.

The recurrence gives intertwining on finite sums, then everywhere
by boundedness. The spectral support is the entire closed interval,
so the operator norm is $2\sqrt p$.
Bounded Borel functional calculus proves the state identity.
\end{proof}

For the Bruhat--Tits tree of $\mathrm{PGL}_2(\mathbb Q_p)$,
vertices are homothety classes of lattices in $\mathbb Q_p^2$.
Adjacency is represented by
$pL\subset L'\subset L$ with $[L:L']=p$.
Lines in $L/pL$ give $p+1$ neighbors.
The maximal compact subgroup has volume one, so this adjacency
is the Hecke operator of $\operatorname{diag}(p,1)$.
Pullback to $\mathrm{GL}_2$ makes $\mathsf z_p$ act as the identity.
The radial transform consequently preserves both the spherical
Hecke action and its root state.

\begin{proposition}[Distance operators]
\label{v4-arith:hecke:prop:distance}
Define $\mathcal D_r$ by
\[
 (\mathcal D_rf)(v)=\sum_{d(v,w)=r}f(w).
\]
Then $\mathcal D_r=R_r(T)$ and
\[
 H_k(T,p)=\sum_{j=0}^{\lfloor k/2\rfloor}\mathcal D_{k-2j}.
\]
Their root expectations satisfy
\[
 \langle e_0,H_k(T,p)e_0\rangle=
 \begin{cases}1,&k\text{ even},\\0,&k\text{ odd}.
 \end{cases}
\]
On $\mathcal R_p$, the generating series converges in operator norm
for $|t|<p^{-1/2}$, and
\[
 \left\langle e_0,(I-tT+pt^2I)^{-1}e_0\right\rangle
 =\frac1{1-t^2}.
\]
\end{proposition}
\begin{proof}
Every row and column of $\mathcal D_r$ has $s_r$ entries,
so it is bounded. Counting the first step gives
\[
 T\mathcal D_1=\mathcal D_2+(p+1)I,\qquad
 T\mathcal D_r=\mathcal D_{r+1}+p\mathcal D_{r-1}\quad(r\geq2).
\]
Hence $\mathcal D_r=R_r(T)$.
With $H_{-1}=H_{-2}=0$, the recurrences give $R_r=H_r-H_{r-2}$.
Telescoping proves the sum formula.
Only $\mathcal D_0e_0$ has a nonzero inner product with $e_0$,
giving the expectations.

The spectral values are
\[
 H_k(2\sqrt p\cos\theta,p)
 =p^{k/2}\sum_{j=0}^ke^{i(k-2j)\theta}.
\]
They include the coinciding-parameter cases $\theta=0,\pi$.
Thus $\|H_k(T,p)\|\leq(k+1)p^{k/2}$ on $\mathcal R_p$.
This proves norm convergence in the stated disk.
Multiplying the series by $I-tT+pt^2I$ proves the inverse formula,
and its state can be evaluated term by term.
\end{proof}

At $p=2$, the first normalized radial polynomials are
\[
 1,\qquad x/\sqrt3,\qquad (x^2-3)/\sqrt6,
\]
whereas $H_2(x,2)=x^2-2$.
There are three two-step returns to the root.
Of the four-step returns, nine return after two steps and six
first reach distance two. Therefore
\[
 \int x^2d\nu_2=3,\qquad \int x^4d\nu_2=15,\qquad
 \int H_2d\nu_2=1,\quad \int R_2d\nu_2=0.
\]
The distinction persists before taking a spectral integral.

There is also a distinction between two cyclic Hecke representations.
Give the free complex mode space the Hermitian form
$\langle x_k,x_l\rangle=k\delta_{kl}$ and complete it.
Then $U_p^*=pB_p$, and $\mathsf H_p=U_p+pB_p$ is bounded and self-adjoint.
At its unit vector $x_1$,
\[
 \langle x_1,\mathsf H_p^2x_1\rangle=p,
 \qquad \langle e_0,T^2e_0\rangle=p+1.
\]
The first value follows from
$\mathsf H_p^2x_1=x_{p^2}+px_1$.
Thus the common Hecke recursion does not give an intertwiner
preserving these two specified states.

\subsection{Finite sets of primes}

For a finite prime set $F$, tensor the radial maps to obtain
\[
 \mathcal U_F:\widehat\bigotimes_{p\in F}\mathcal R_p
 \xrightarrow{\sim}
 L^2\left([0,\pi]^F,\prod_{p\in F}\nu_p\right).
\]
It sends $\bigotimes_pe_{r_p}$ to the product of the normalized
radial polynomials. It preserves each local Hecke operator and
sends the tensor root state to integration against the product measure.
The integral cyclic lattices are the finite tensor products of
the sphere-sum lattices.

Adjoining a prime uses $v\mapsto v\otimes e_0$ on the tree side
and $f\mapsto f\otimes1$ on the spectral side.
These isometries commute with $\mathcal U_F$ and preserve the states.
Their algebraic union therefore has an isometric comparison,
which extends to its Hilbert completion.
Each finite-support Hecke operation acts on its own natural factors.
No bound uniform over all primes is required or asserted.

\section{The integral symmetric-function lattice}
\label{v4-arith:hecke:sec:lambda}

The index shifts already form a monoid action on $P_R$.
Their integral Frobenius property depends on the lattice.
For one species over $\mathbb Z$,
\[
 \psi_px_1-x_1^p=x_p-x_1^p\notin pP_{\mathbb Z}.
\]
The two monomials remain distinct after reduction modulo $p$.
We now construct the larger lattice on which the same rational
operations satisfy the required congruences.

Inside $P_{\mathbb Q}$, define $h_0=1$ and
\begin{equation}
 \label{v4-arith:hecke:eq:complete-series}
 H(t)=\sum_{r\geq0}h_rt^r
 =\exp\left(\sum_{k\geq1}\frac{x_kt^k}{k}\right),\qquad
 \Lambda=\mathbb Z[h_1,h_2,\ldots].
\end{equation}
The coefficients are defined formally, one degree at a time.
Differentiation gives Newton's recursion
\[
 rh_r=\sum_{k=1}^rx_kh_{r-k}.
\]
It expresses $x_r$ as an integral polynomial in $h_1,\ldots,h_r$.
Conversely, $h_r=x_r/r$ plus a polynomial in earlier $x_k$ over $\mathbb Q$.

Thus the $h_r$ are algebraically independent and
\[
 P_{\mathbb Z}\subsetneq\Lambda,\qquad
 P_{\mathbb Q}=\Lambda\otimes\mathbb Q.
\]
For instance, $h_2=(x_1^2+x_2)/2$ does not belong to $P_{\mathbb Z}$.
We call $\Lambda$ the \emph{Hall lattice}.
For several species, take its finite tensor products.

To identify its standard operations, specialize
$x_k=z_1^k+\cdots+z_N^k$.
Then $H(t)=\prod_{i=1}^N(1-z_it)^{-1}$, so $h_r$ is the complete
symmetric polynomial of degree $r$.
These symmetric polynomials stabilize when zero variables are adjoined.
Every fixed weight is detected once $N$ is sufficiently large.
Substitution $z_i\mapsto z_i^n$ therefore preserves $\Lambda$
and induces $\psi_nx_k=x_{nk}$.

In characteristic $p$, the same substitution is the Frobenius map
on integer symmetric polynomials modulo $p$. Hence
\begin{equation}
 \label{v4-arith:hecke:eq:frobenius}
 \psi_p(f)\equiv f^p\pmod{p\Lambda}\qquad(f\in\Lambda).
\end{equation}

\subsection{Special lambda operations}

The full integral structure includes operations beyond their rational
power sums. We specify the convention explicitly.
For two finite alphabets $X=(z_i)$ and $Y=(w_j)$, let $e_n(X)$
denote the elementary symmetric polynomial.
There are universal integer polynomials $P_n$ and $Q_{m,n}$ determined by
\[
 e_n(\{z_iw_j\}_{i,j})
 =P_n(e_1(X),\ldots,e_n(X);e_1(Y),\ldots,e_n(Y)),
\]
and
\[
 e_m(\{z_{i_1}\cdots z_{i_n}:i_1<\cdots<i_n\})
 =Q_{m,n}(e_1(X),\ldots,e_{mn}(X)).
\]
The fundamental theorem of symmetric polynomials gives their integral
coefficients and stability as the alphabets grow.

\begin{definition}
\label{v4-arith:hecke:def:lambda}
A \emph{special $\lambda$-ring} is a commutative unital ring $A$
with operations $\lambda^n:A\to A$ for $n\geq0$.
Writing $\lambda_t(a)=\sum_{n\geq0}\lambda^n(a)t^n$, their identities are
\[
 \lambda^0(a)=1,\qquad \lambda^1(a)=a,\qquad
 \lambda_t(a+b)=\lambda_t(a)\lambda_t(b),\qquad
 \lambda_t(1)=1+t,
\]
together with the product and composition identities given by
$P_n$ and $Q_{m,n}$ above.
Its \emph{Adams operations} are the integral coefficient operations in
\[
 -t\frac{(\lambda_{-t}(a))'}{\lambda_{-t}(a)}
 =\sum_{n\geq1}\psi_n(a)t^n.
\]
Here the prime means formal differentiation in $t$.
\end{definition}

The alphabet identities show that each Adams operation is a unital
ring endomorphism and $\psi_m\psi_n=\psi_{mn}$.
They also give the Frobenius congruences.
These are identities of universal integer polynomials, so they remain
valid after evaluation in every special $\lambda$-ring.
The ordinary symmetric functions give such a structure on $\Lambda$.
Their Adams operations are precisely the index shifts above.

\begin{theorem}[Torsion-free Adams criterion]
\label{v4-arith:hecke:thm:wilkerson}
Let the additive group of $A$ be torsion-free.
A family of unital ring endomorphisms $\psi_n$ determines a unique
special $\lambda$-structure if
\[
 \psi_1=I,\qquad \psi_m\psi_n=\psi_{mn},\qquad
 \psi_p(a)\equiv a^p\pmod{pA}\quad\text{for every prime }p.
\]
\end{theorem}

This is Wilkerson's criterion, as stated in Borger,
\emph{The basic geometry of Witt vectors, I}, Section~1.17.
Over a $\mathbb Q$-algebra its reconstruction takes the form
\[
 \lambda_t(a)=\exp\left(
 \sum_{n\geq1}\frac{(-1)^{n-1}\psi_n(a)t^n}{n}\right).
\]
Over $\mathbb Z$, the shifts therefore give the special structure
on $\Lambda$, but fail the necessary congruence on $P_{\mathbb Z}$.
This failure concerns the specified shifts.
For example, the different maps $x_k\mapsto x_k^n$ give
commuting Frobenius lifts on $P_{\mathbb Z}$.
They define the structure in which each variable is a line element,
meaning $\lambda_t(x_k)=1+x_kt$.

\begin{example}[The torsion hypothesis]
\label{v4-arith:hecke:ex:torsion-lambda}
On $\mathbb F_p$, let every $\psi_n$ be the identity.
These ring maps commute and satisfy all the displayed Frobenius congruences.
At $p$, this uses $a^p=a$. At other primes $q$, the ideal
$q\mathbb F_p$ is the whole ring.

They cannot come from a special $\lambda$-structure.
Such a structure would give
\[
 1=\lambda_t(0)=\lambda_t(p\cdot1)
   =\lambda_t(1)^p=(1+t)^p=1+t^p,
\]
which is false in $\mathbb F_p[[t]]$.
Thus commuting Frobenius lifts alone do not replace the full
integral operations on a ring with additive torsion.
\end{example}

\begin{proposition}[Divisor substitutions on the Hall lattice]
\label{v4-arith:hecke:prop:hall-lattice}
Every $D_n$ preserves $\Lambda$, and
\begin{equation}
 \label{v4-arith:hecke:eq:divisor-series}
 D_nH(t)=\prod_{d\mid n}\psi_nH(t^d).
\end{equation}
Nevertheless $D_p$ is not a Frobenius lift on this lattice:
\[
 D_ph_p=\psi_ph_p+\psi_ph_1
 \equiv h_p^p+h_1^p\not\equiv h_p^p\pmod{p\Lambda}.
\]
\end{proposition}
\begin{proof}
In $P_{\mathbb Q}[[t]]$, set $k=dj$ in the logarithm to obtain
\[
 \sum_{k\geq1}\frac{D_nx_k}{k}t^k
 =\sum_{d\mid n}\sum_{j\geq1}\frac{x_{nj}}j t^{dj}.
\]
Exponentiation gives \eqref{v4-arith:hecke:eq:divisor-series}.
Every factor has integral Hall coefficients, proving preservation of $\Lambda$.
For $n=p$, only the two stated contributions have degree $p$.
Apply \eqref{v4-arith:hecke:eq:frobenius} to them.
The term $h_1^p$ is nonzero in the polynomial ring $\Lambda/p\Lambda$.
\end{proof}

In particular, the equality $D_{p^a}=\psi_{p^a}$ on
$P_{\mathbb F_p}$ does not imply equality after reducing $\Lambda$.
The inclusion $P_{\mathbb Z}\subset\Lambda$ is not an identification
of their integral coordinates or reductions.

There is a useful comparison which does preserve the algebraic images.
Define a triangular substitution on $P_R$ by
\[
 C_nx_{s,k}=\sum_{d\mid(n,k)}d\,x_{s,k/d}.
\]
Its leading term is $x_{s,k}$, and every other mode index is smaller.
Induction therefore constructs an integral polynomial inverse on
every generator, over every coefficient ring.

The same logarithmic calculation gives
\[
 C_nH(t)=\prod_{d\mid n}H(t^d).
\]
Its coefficient of $t^r$ is $h_r$ plus an integral polynomial
in $h_1,\ldots,h_{r-1}$.
Thus $C_n$ is also an automorphism of the Hall lattice.

Direct evaluation on the generators now proves
\begin{equation}
 \label{v4-arith:hecke:eq:triangular}
 D_n=\psi_nC_n,\qquad
 D_n(P_R)=\psi_n(P_R),\qquad
 D_n(\Lambda)=\psi_n(\Lambda).
\end{equation}
Equal images do not identify the two actions or their composition laws.
Nor does the triangular inverse have to survive a Hilbert completion.

\section{Big Witt coordinates and the cofree construction}
\label{v4-arith:hecke:sec:witt}

The previous exponential formulas live first over $\mathbb Q$.
Big Witt coordinates give integral formulas before any scalar extension.
For a commutative ring $R$, put $W(R)=\prod_{r\geq1}R$ as a set.
Associate to $a=(a_r)$ the series
\begin{equation}
 \label{v4-arith:hecke:eq:witt-series}
 A_a(t)=\prod_{r\geq1}(1-a_rt^r)^{-1}\in1+tR[[t]].
\end{equation}

Each coefficient involves finitely many factors.
Its coefficient of $t^r$ is $a_r$ plus a polynomial in earlier coordinates.
Induction therefore gives a natural bijection
$W(R)\simeq1+tR[[t]]$ over every ring.

The \emph{ghost components} are
\begin{equation}
 \label{v4-arith:hecke:eq:ghosts}
 w_n(a)=\sum_{r\mid n}r\,a_r^{n/r},\qquad
 t\frac{A_a'(t)}{A_a(t)}=\sum_{n\geq1}w_n(a)t^n.
\end{equation}
The derivative identity follows by differentiating the factors.
It involves no division by an integer in $R$.

\begin{proposition}
\label{v4-arith:hecke:prop:ghosts}
If the additive group of $R$ is torsion-free, the ghost map is injective.
Over $\mathbb F_p$, the nonzero vector with $a_p=1$ and all other
coordinates zero has every ghost equal to zero.
Its series is $(1-t^p)^{-1}\ne1$.
\end{proposition}
\begin{proof}
If two vectors agree below $n$, equality of their $n$th ghosts gives
$n(a_n-b_n)=0$.
Torsion-freeness gives equality at $n$, and induction proves injectivity.
For the stated vector over $\mathbb F_p$, the only possible ghost
contribution is $p$ when $p\mid n$.
It vanishes, although the coefficient of $t^p$ in the series is one.
\end{proof}

We construct the ring operations on the series themselves.
Write $E_r(a)=(1-at^r)^{-1}$.
Witt addition is ordinary multiplication of series:
$A+_WB=AB$. Its zero is the series $1$, and its negative is $A^{-1}$.

Witt multiplication, denoted $\odot$, is determined by
\begin{equation}
 \label{v4-arith:hecke:eq:witt-product}
 E_r(a)\odot E_s(b)
 =\left(1-a^{s/g}b^{r/g}t^l\right)^{-g},
 \qquad g=(r,s),\quad l=\operatorname{lcm}(r,s),
\end{equation}
and distributivity over Witt addition.
For two arbitrary coordinate products, multiply these factors
over all pairs $(r,s)$.
Since $l\geq r,s$, each coefficient uses finitely many pairs.
The construction is therefore a universal integral polynomial operation.

\begin{theorem}[The big Witt ring]
\label{v4-arith:hecke:thm:witt-ring}
These operations make $W(R)$ a natural commutative unital ring.
Its multiplicative unit is $E_1(1)$.
Every $w_n:W(R)\to R$ is a ring homomorphism.
\end{theorem}
\begin{proof}
Logarithmic differentiation makes each ghost additive under series multiplication.
For one factor,
\[
 w_n(E_r(a))=
 \begin{cases}r a^{n/r},&r\mid n,\\0,&r\nmid n.\end{cases}
\]
The ghost of \eqref{v4-arith:hecke:eq:witt-product} is
$gl\,a^{n/r}b^{n/s}$ when $l\mid n$, and zero otherwise.
Since $gl=rs$, this is the product of the two ghosts.
Distributivity then proves multiplicativity for arbitrary series.

It remains to verify the ring identities without assuming injectivity
on the given coefficient ring.
Fix a coefficient degree and regard the finitely many input coordinates
as independent variables over $\mathbb Z$.
This polynomial ring is torsion-free, so its ghosts are injective.
Commutativity, associativity, distributivity and the unit law hold on ghosts,
hence on these universal coordinate polynomials.
Specializing their integer coefficients proves the identities over every $R$.
\end{proof}

The Teichm\"uller element is $[a]=E_1(a)$.
It satisfies $[a]\odot[b]=[ab]$, but need not preserve addition in $R$.
For $m\geq1$, define Verschiebung by $V_mA(t)=A(t^m)$.
It is additive for $+_W$.
Define Frobenius by its values on the coordinate factors:
\begin{equation}
 \label{v4-arith:hecke:eq:witt-frobenius}
 F_mE_r(a)=
 \left(1-a^{m/g}t^{r/g}\right)^{-g},
 \qquad g=(m,r),
\end{equation}
and extend by products of series.
A coefficient through degree $N$ uses only $r\leq mN$,
so this also defines an integral operation coefficientwise.

\begin{proposition}[Frobenius and Verschiebung]
\label{v4-arith:hecke:prop:fv}
Each $F_m$ is a natural unital ring endomorphism of $W(R)$, and
\[
 w_n(F_ma)=w_{mn}(a),\qquad
 w_n(V_ma)=
 \begin{cases}m\,w_{n/m}(a),&m\mid n,\\0,&m\nmid n.\end{cases}
\]
Their identities over every coefficient ring include
\begin{align*}
 F_mF_n&=F_{mn},& V_mV_n&=V_{mn},& F_mV_m&=mI,\\
 F_mV_n&=gV_{n/g}F_{m/g},&&g=(m,n).
\end{align*}
The integer multiple means repeated Witt addition.
In particular, $F_m[a]=[a^m]$.
\end{proposition}
\begin{proof}
For a single factor, \eqref{v4-arith:hecke:eq:witt-frobenius}
has $n$th ghost $r a^{mn/r}$ precisely when $r\mid mn$.
This proves the Frobenius formula.
Substituting $t^m$ in the logarithmic derivative proves the Verschiebung formula.
The ring-map property and all displayed identities follow on ghosts.
Each operation was already constructed by universal integer polynomials.
The torsion-free universal-ring argument from the theorem therefore
proves the identities before specialization to arbitrary $R$.
\end{proof}

\subsection{Plethysm and the full comonad}

The ring $\Lambda=\mathbb Z[h_1,h_2,\ldots]$ represents the series functor:
\[
 W(R)\simeq\operatorname{Hom}_{\mathrm{Ring}}(\Lambda,R),\qquad
 f\longmapsto\sum_{n\geq0}f(h_n)t^n.
\]
The values $f(h_n)$ are series coefficients, not the Witt coordinates $a_n$.
They are related by the triangular integral bijection
\eqref{v4-arith:hecke:eq:witt-series}.
The power sum $x_n\in\Lambda$ evaluates to $w_n$.

We need the integral composition operation on $\Lambda$.
For $h\in\Lambda$, write its specialization to a finite alphabet
as an integer sum of monomials, $h(X)=\sum_Mc_MM$.
Define
\[
 \sum_{n\geq0}(h_n\circ h)(X)t^n
 =\prod_M(1-Mt)^{-c_M}.
\]
Negative multiplicities are allowed, using the inverse formal series.
The coefficients are integral symmetric polynomials and stabilize
when zero variables are added.

Extend $g\mapsto g\circ h$ as a ring homomorphism from $\Lambda$.
This operation is \emph{plethysm}.
Its unit is $h_1=x_1$.

Plethysm is associative. One way to check this precisely is to
embed $\Lambda$ in $\mathbb Q[x_1,x_2,\ldots]$.
The defining formula gives $x_n\circ h=\psi_n(h)$,
where $\psi_nx_k=x_{nk}$ and rational coefficients are fixed.
These substitutions commute and satisfy
$\psi_n(g\circ h)=(\psi_ng)\circ h$.
Evaluating a polynomial in the $x_n$ therefore proves
$(f\circ g)\circ h=f\circ(g\circ h)$.
Since the operation was constructed integrally, the equality lies in $\Lambda$.

\begin{theorem}[Cofree special lambda rings]
\label{v4-arith:hecke:thm:cofree}
The functor $W$ is right adjoint to the forgetful functor
from special $\lambda$-rings to commutative rings.
Its counit and comultiplication are
\begin{equation}
 \label{v4-arith:hecke:eq:comonad}
 \varepsilon_R(f)=f(h_1),\qquad
 (\Delta_Rf)(g)(h)=f(h\circ g).
\end{equation}
Thus $\Delta_R:W(R)\to W(W(R))$.
It satisfies both counit identities and
$\Delta_{W(R)}\Delta_R=W(\Delta_R)\Delta_R$.
Its outer ghosts are $w_n(\Delta_Rf)=F_nf$ as elements of $W(R)$.
\end{theorem}
\begin{proof}
For $f\in W(R)$ and $g\in\Lambda$, define the operation
$g(f)\in W(R)$ by $g(f)(h)=f(h\circ g)$.
For fixed $g$, the right side is a ring homomorphism in $h$,
by the definition of plethysm.
The universal symmetric-polynomial identities for sums and products
show that $g\mapsto g(f)$ is a ring homomorphism into $W(R)$.
These identities are exactly the series addition and multiplication
constructed above. Associativity of plethysm supplies composition
of the operations. Taking the elementary symmetric operations
therefore makes $W(R)$ a special $\lambda$-ring.

Let $A$ be a special $\lambda$-ring and $\phi:A\to R$ a ring map.
For $a\in A$, evaluation of symmetric operations gives the map
\[
 \widehat\phi(a)(g)=\phi(g(a))\qquad(g\in\Lambda).
\]
The sum and product identities make
$\widehat\phi:A\to W(R)$ a ring map.
The composition identities make it a $\lambda$-ring map.
Its counit is $\phi$, since $h_1(a)=a$.

Conversely, the counit determines a $\lambda$-ring map into $W(R)$.
Applying it after the operation $g$ recovers evaluation at $g$.
This proves the natural adjunction.

Applying this construction to the identity on the underlying ring
of $W(R)$ gives the stated $\Delta_R$.
The equalities $h_1\circ g=g$ and $h\circ h_1=h$ prove
the two counit identities separately.
Evaluating the two coassociativity maps on three elements of $\Lambda$
reduces them to the two associations of the same plethysm.
Finally, $h\circ x_n=\psi_n(h)$.
Evaluation of the outer power sum $x_n$ in \eqref{v4-arith:hecke:eq:comonad}
therefore gives precisely the coefficient action $F_n$.
\end{proof}

Equivalently, a special $\lambda$-structure on a ring $A$ is a ring map
$\eta:A\to W(A)$ satisfying
\[
 \varepsilon_A\eta=I,
 \qquad \Delta_A\eta=W(\eta)\eta.
\]
Indeed, its coordinates define the complete symmetric operations.
The fact that $\eta$ is a ring map gives their sum and product identities.
The second displayed equality gives composition, and the first fixes
the degree-one operation. Inverting the complete symmetric series
at $-t$ recovers the elementary operations of
Definition~\ref{v4-arith:hecke:def:lambda}.
The converse follows by evaluation of those operations.

The two universal constructions have different variances:
\[
 \operatorname{Hom}_{\lambda\text{-}\mathrm{Ring}}(\Lambda,A)
 \simeq A,
 \qquad
 \operatorname{Hom}_{\lambda\text{-}\mathrm{Ring}}(A,W(R))
 \simeq\operatorname{Hom}_{\mathrm{Ring}}(A,R).
\]
The first says that $\Lambda$ is free on one distinguished element.
The second is the cofree property of $W$.
On rings with additive torsion, the full map $\eta$ and its
composition identity retain data which its ghost operations can lose.
Borger's descent interpretation uses this integral structure and its
forgetful functor. It does not create an arithmetic comparison from
a rational index-shift formula.

\section{Products, gradings, and algebraic completions}
\label{v4-arith:hecke:sec:structures}

The polynomial ring has a primitive Hopf structure over $R$:
\[
 \Delta x_{s,k}=x_{s,k}\otimes1+1\otimes x_{s,k},\qquad
 \epsilon(x_{s,k})=0,\qquad \mathcal S(x_{s,k})=-x_{s,k}.
\]
Here $\mathcal S$ denotes the antipode.
Each of $D_n$, $\psi_n$ and $C_n$ preserves this structure.
Their generator images are primitive, and the identities of algebra
maps can be checked on generators.
They also preserve the unit, augmentation and particle degree.

On the Hall lattice, the same coproduct is
\[
 \Delta h_r=\sum_{i+j=r}h_i\otimes h_j.
\]
Preservation follows from the rational power-sum identity and the
torsion-free integral inclusion.
This is the ordinary symmetric-function Hopf structure.
An arithmetic extension correspondence would require its own comparison.

\begin{proposition}[Weights and completions]
\label{v4-arith:hecke:prop:weights}
For a nonempty species set and any nonzero coefficient ring,
$D_n$ and $\psi_n$ are injective on $P_R$.
For $n>1$, neither is surjective.
The index shift sends weight $w$ to weight $nw$.
The divisor substitution sends weight $w$ into weights between $w$ and $nw$.
Both act on the finite-weight quotients and continuously on
\[
 \widehat P_R^{\mathrm{wt}}=\prod_{w\geq0}P_{R,w}.
\]
They also extend to the coefficient-adic completion
$\varprojlim_jP_{\mathbb Z}/\ell^jP_{\mathbb Z}$ for every prime $\ell$.
\end{proposition}
\begin{proof}
Let $f$ have highest nonzero weight $w$.
The highest-weight term of $D_nf$ is $\psi_n$ applied to the
weight-$w$ component, and it has weight $nw$.
Index substitution is injective on the monomial basis, so this
component cannot vanish over a nonzero coefficient ring.
This proves injectivity of both maps.
For $n>1$, the same highest-weight argument excludes a preimage of $x_1$.

Every term of $D_nx_k$ has weight between $k$ and $nk$.
Multiplication proves the weight bounds on monomials.
Thus the ideal of weights greater than $W$ is preserved.
Its quotient has finite rank for a finite species set.
An output through weight $W$ depends only on inputs through weight $W$,
which proves continuity on the weight completion.

Finally, the integral maps preserve every ideal $\ell^jP_{\mathbb Z}$.
Their compatible quotient maps give the coefficient-adic extensions.
\end{proof}

A finite subalgebra generated by modes $k\leq N$ is not invariant
when $n>1$: the highest term of $D_nx_N$ is $x_{nN}$.
The finite-weight quotient is a different construction.
The coefficient-adic completion is also different from the weight product.
In monomial coordinates, its coefficients lie in $\mathbb Z_\ell$,
with only finitely many nonzero coefficients modulo each $\ell^j$.
The unrestricted weight product imposes no such decay.

Other symmetric-function operations give different preservation laws.
For one species over $\mathbb Q$, the internal coproduct is the
algebra map $\delta x_k=x_k\otimes x_k$.
It is preserved by $\psi_n$.
For $D_n$ with $n>1$, the image of $x_n$ contains both
$x_{n^2}$ and $nx_n$.
The tensor square of that image has cross terms absent from
$\delta(D_nx_n)$, so $D_n$ does not preserve $\delta$.

For a partition $\lambda$ with multiplicities $m_k$, write
$x_\lambda=\prod_kx_k^{m_k}$ and $z_\lambda=\prod_kk^{m_k}m_k!$.
The internal product is the bilinear operation
\[
 x_\lambda*x_\mu=\delta_{\lambda\mu}z_\lambda x_\lambda.
\]
Neither family preserves it for $n>1$.
Indeed, $x_1*x_1=x_1$, whereas $x_n*x_n=nx_n$.
These nonvanishing obstructions use characteristic zero.
They do not transfer unchanged to a torsion coefficient ring.

\section{Oscillator intertwiners and vertex products}
\label{v4-arith:hecke:sec:oscillators}

For one species, define polynomial operators
\[
 a_{-k}f=x_kf,\qquad a_kf=k\partial_{x_k}f\quad(k\geq1),
 \qquad a_0=0.
\]
They obey $[a_k,a_{-l}]=k\delta_{kl}$, with central element acting as one.
Different species commute.
The Hall pairing is
$\langle x_\lambda,x_\mu\rangle=\delta_{\lambda\mu}z_\lambda$.
Over $\mathbb C$, its Hermitian version is positive and makes
$a_k$ the adjoint of $a_{-k}$ on polynomials.

The substitutions intertwine creation by $x_k$ with creation
by its image. Their annihilation intertwiners require a separate calculation.
For the index shifts, differentiation gives
\begin{equation}
 \label{v4-arith:hecke:eq:adams-oscillators}
 \psi_na_{-k}=a_{-nk}\psi_n,\qquad
 a_{nk}\psi_n=n\psi_na_k,\qquad
 a_j\psi_n=0\quad(n\nmid j).
\end{equation}
Sending both positive and negative modes to their index multiples
would multiply the commutator by $n$.
Over $\mathbb Q$, the assignments
\[
 a_{-k}\longmapsto a_{-nk},\qquad
 a_k\longmapsto a_{nk}/n
\]
instead define a central-charge-one Heisenberg embedding.
It is intertwined by $\psi_n$, but does not preserve the Hall adjoint.

\begin{proposition}[An integral annihilation lift]
\label{v4-arith:hecke:prop:annihilation}
Fix $n\geq1$. Define integers $c_{k,i}$ in increasing order of $i$ by
\[
 c_{k,i}=0\quad(n\nmid i),\qquad
 c_{k,nj}=k\delta_{kj}
 -\sum_{\substack{d\mid(n,j)\\d>1}}d\,c_{k,nj/d}.
\]
The polynomial derivations
\[
 \mathcal A_k^{(n)}=\sum_{i\geq1}c_{k,i}\partial_{x_i}
\]
are well-defined and commute. They satisfy
\[
 \mathcal A_k^{(n)}(D_nx_j)=k\delta_{kj},\qquad
 \mathcal A_k^{(n)}D_n=D_na_k.
\]
Together with multiplication by $D_nx_j$, they give
a central-charge-one oscillator representation.
\end{proposition}
\begin{proof}
Every term on the right of the recursion has smaller index.
It therefore defines all coefficients integrally.
Each polynomial uses finitely many variables, so only finitely
many derivatives contribute to its image.
Constant-coefficient derivations commute.

Substituting \eqref{v4-arith:hecke:eq:substitutions} into the first identity
gives exactly the recursion.
The chain rule proves the second identity on all polynomials.
Commuting a derivation past multiplication by $D_nx_j$ gives
its constant derivative $k\delta_{kj}$, proving the oscillator relations.
\end{proof}

Over $\mathbb Q$, the lift can be written
$\mathcal A_k^{(n)}=\sum_i(c_{k,i}/i)a_i$.
Integrality concerns the polynomial derivation, not these coefficients
in the fixed annihilation modes.
The lift is generally not adjoint-preserving.
Both unscaled substitutions send $x_1$ to $x_n$, whose squared norm is $n$.
Creation by $x_n$ and its Hall adjoint therefore have commutator $n$,
rather than one. This obstructs the specified adjoint-preserving
assignment when $n>1$.

\begin{proposition}[The vertex obstruction]
\label{v4-arith:hecke:prop:vertex}
Over $\mathbb Q$ or $\mathbb C$, neither $D_n$ nor $\psi_n$
is an endomorphism of the usual Heisenberg vertex algebra for $n>1$.
Both remain polynomial algebra and primitive Hopf endomorphisms.
\end{proposition}
\begin{proof}
The standard translation operator satisfies $\partial x_k=kx_{k+1}$.
Every vertex homomorphism preserves it, since
$\partial v=v_{(-2)}1$.
Here
\[
 \partial\psi_nx_1=nx_{n+1},\qquad
 \psi_n\partial x_1=x_{2n},
\]
and
\[
 \partial D_nx_1=nx_{n+1},\qquad
 D_n\partial x_1=x_{2n}+2\mathbf1_{2\mid n}x_n.
\]
The values differ for every $n>1$.
This argument does not require preservation of a chosen conformal vector.

The same failure is visible in the singular products.
Writing $a(z)=Y(x_1,z)$, one has
\[
 Y(x_n,z)=\frac{\partial_z^{n-1}a(z)}{(n-1)!},\qquad
 Y(x_n,z)Y(x_n,w)\sim
 \frac{(-1)^{n-1}(2n-1)!}
 {((n-1)!)^2(z-w)^{2n}}.
\]
Differentiate $a(z)a(w)\sim(z-w)^{-2}$ to obtain this coefficient.
For $n>1$, the higher pole cannot be the image of the original
central second-order pole under a vacuum-preserving vertex map.
The polynomial and Hopf assertions were proved above.
\end{proof}

For a finite species set, the index shifts nevertheless give a
vacuum-preserving action of $\mathbb N^\times$ on the underlying
vector space. This means endomorphisms with $\psi_1=I$ and
$\psi_m\psi_n=\psi_{mn}$.
It does not assert an action by vertex endomorphisms.
The divisor family fails even that monoid law in characteristic zero.

\section{Hilbert closures and a weighted threshold}
\label{v4-arith:hecke:sec:hilbert}

Fix a nonempty finite species set and complete the complex polynomial
space in the Hermitian Hall form.
It is the symmetric Hilbert Fock space
$\mathcal F=\bigoplus_{j\geq0}\mathcal F^{(j)}$,
with squared monomial norms $\prod_{s,k}k^{m_{s,k}}m_{s,k}!$.
Its one-particle space has $\|x_{s,k}\|^2=k$.
The map $\psi_n/\sqrt n$ is an isometry there.

Each divisor summand
$x_k\mapsto d\mathbf1_{d\mid k}x_{nk/d}$ has norm $\sqrt{nd}$.
Its nonzero basis images are orthogonal with that constant norm ratio.
Consequently, for $d_n=D_n|_{M_{\mathbb C}}$,
\[
 \|d_n\|\leq\sum_{d\mid n}\sqrt{nd}.
\]
Both maps are bounded on one particle and on every fixed
particle sector after taking symmetric tensor powers.

\begin{proposition}[Closed full-Fock operators]
\label{v4-arith:hecke:prop:fock-closure}
For $A=d_n$ or $A=\psi_n|_{M_{\mathbb C}}$, the polynomial operator
$\operatorname{Sym}(A)$ is closable.
Its closure is $\Gamma(A)=\bigoplus_{j\geq0}\operatorname{Sym}^j(A)$,
with domain
\[
 \operatorname{Dom}\Gamma(A)=
 \left\{(v_j)\in\mathcal F:
       \sum_{j\geq0}\|\operatorname{Sym}^j(A)v_j\|^2<\infty\right\}.
\]
For $n>1$, both closures are unbounded.
For the index shift, the domain condition is exactly
$\sum_jn^j\|v_j\|^2<\infty$.
\end{proposition}
\begin{proof}
A direct sum of bounded sector maps on the displayed domain is closed.
Finite particle truncation approximates a vector and its image.
Within each retained sector, finite-mode polynomials approximate
both, since the sector map is bounded.
Thus polynomials are a core and the direct sum is precisely their closure.

For the index shift, each sector map multiplies norms by $n^{j/2}$.
Both substitutions send $x_1^j$ to $x_n^j$, giving
\[
 \frac{\|D_nx_1^j\|}{\|x_1^j\|}
 =\frac{\|\psi_nx_1^j\|}{\|x_1^j\|}=n^{j/2}.
\]
This proves unboundedness for $n>1$ and the stated domain formula.
\end{proof}

The normalized substitution
$\widetilde\psi_nx_{s,k}=n^{-1/2}x_{s,nk}$ gives a different positive result.
It is an isometry on every particle sector and their Hilbert direct sum.
It is also a complex polynomial Hopf map and a monoid action.
Sending both $a_{-k}$ and $a_k$ to their index multiples divided
by $\sqrt n$ preserves the Heisenberg relations and Hall adjoints.
Equation~\eqref{v4-arith:hecke:eq:adams-oscillators}, with this normalization,
shows that $\widetilde\psi_n$ intertwines the resulting embedding.
The normalization changes the integral structure and still fails
ordinary translation for $n>1$.

\subsection{The kernel created by completion}

The triangular automorphism $C_n$ in
\eqref{v4-arith:hecke:eq:triangular} has a bounded one-particle extension.
For $B_dx_k=\mathbf1_{d\mid k}x_{k/d}$, one has
\[
 c_n=C_n|_M=\sum_{d\mid n}dB_d,
 \qquad \|c_n\|\leq\sum_{d\mid n}\sqrt d,
 \qquad c_n=\prod_{p\mid n}c_{p^{v_p(n)}}.
\]
The norm estimate uses $\|B_d\|=d^{-1/2}$.
The factorization follows from $B_dB_e=B_{de}$.

\begin{theorem}[One-particle kernel and range]
\label{v4-arith:hecke:thm:hilbert-kernel}
The operator $c_n$ is a bounded surjection on the completed
one-particle space. The divisor operator has closed range
\[
 \operatorname{Ran}d_n
 =\overline{\operatorname{span}}\{x_{nk}:k\geq1\}
 =\operatorname{Ran}(\psi_n|_M),\qquad
 \ker d_n=\ker c_n.
\]
For $n>1$, this kernel is infinite-dimensional on the full
one-species space.
On a prime chain $e_r=x_{up^r}$ and for $a\geq1$,
\[
 c_{p^a}=\sum_{j=0}^ap^jB^j,
 \qquad
 \ker c_{p^a}=\operatorname{span}_{\mathbb C}
 \left\{\sum_{r\geq0}(\zeta/p)^re_r:
       \zeta^{a+1}=1,\ \zeta\ne1\right\}.
\]
The prime-chain kernel has dimension $a$.
For $n=1$, the operators are the identity and the kernels vanish.
\end{theorem}
\begin{proof}
On the chain, $\|e_r\|^2=up^r$ and $\|B\|=p^{-1/2}$.
For each stated root, the geometric vector has squared norm
$u\sum_{r\geq0}p^{-r}=u/(1-p^{-1})$.
It is a $B$-eigenvector with eigenvalue $\zeta/p$,
so $c_{p^a}$ kills it.
Conversely, a kernel vector $\sum_rb_re_r$ satisfies
\[
 \sum_{j=0}^ap^jb_{r+j}=0\qquad(r\geq0).
\]
Its first $a$ coefficients determine all the others.
The distinct roots $\zeta/p$ give independent solutions by the
Vandermonde determinant, exhausting this solution space.

Let $Ue_r=e_{r+1}$, so $BU=I$ and $\|U\|=\sqrt p$.
For $\lambda=\zeta/p$, the operator
\[
 R_\lambda=U(I-\lambda U)^{-1}
\]
is a bounded right inverse of $B-\lambda$.
Indeed, $(B-\lambda)U=I-\lambda U$ and
$\|\lambda U\|=p^{-1/2}<1$.
The inverse is a norm-convergent Neumann series.

Factor
\[
 c_{p^a}=p^a\prod_{\substack{\zeta^{a+1}=1\\\zeta\ne1}}(B-\zeta/p).
\]
Multiplying the right inverses in reverse order, with factor $p^{-a}$,
gives a bounded right inverse of $c_{p^a}$.
Its norm bound is independent of $u$.
It therefore extends over the orthogonal sum of all prime chains.
The finite product over prime factors proves surjectivity of $c_n$.

The factorization $d_n=\psi_nc_n$ and the scaled isometry
$\psi_n|_M$ prove the range and kernel assertions.
For $n>1$, choose one nontrivial prime-power factor.
Its kernel vectors on infinitely many disjoint prime chains are independent.
The other factors commute with it, so all these vectors lie in $\ker c_n$.
\end{proof}

The polynomial inverse has a different domain.
For example,
\[
 C_p^{-1}e_r=\sum_{j=0}^r(-p)^je_{r-j},\qquad
 \frac{\|C_p^{-1}e_r\|^2}{\|e_r\|^2}=\sum_{j=0}^rp^j.
\]
Cancellation in $(I+pB)\sum_{j=0}^r(-pB)^je_r=e_r$
proves the first formula. Orthogonality proves the second.

More generally, the polynomial inverse of $C_n$ is not closable
on this Hilbert space when $n>1$.
Choose $0\ne v\in\ker c_n$ and finite truncations $v_N\to v$.
Then $y_N=c_nv_N\to0$, while $C_n^{-1}y_N=v_N\to v\ne0$.
The graph closure contains $(0,v)$ and cannot be an operator graph.
The same argument applies to the algebraic inverse of $d_n$
on its polynomial image.

These one-particle kernels also lie in $\ker\Gamma(d_n)$.
By contrast, $\Gamma(\psi_n|_M)$ is injective, since each
particle sector is a nonzero scalar multiple of an isometry.
The polynomial maps remain injective throughout.
No extension of polynomial multiplication to all Hilbert vectors is asserted.

\subsection{The exact weighted boundary}

Retain a prime chain $e_r=x_{up^r}$ and its finite linear span $D$.
For $s\geq0$, define
\begin{equation}
 \label{v4-arith:hecke:eq:weighted}
 \mathcal H_s=\left\{\sum_{r\geq0}b_re_r:
   \sum_{r\geq0}(up^r)^{1+2s}|b_r|^2<\infty\right\}.
\end{equation}
If $Ee_r=up^re_r$, this is $\operatorname{Dom}(E^s)$ in $\mathcal H_0$,
with norm $\|E^sv\|_0$.
Since $E\geq1$, this norm is equivalent to its graph norm.
Write $\mathcal H_{s,\geq a}$ for the closed tail spanned by $e_r$, $r\geq a$.
Its codimension is $a$.

\begin{theorem}[Weighted inverses]
\label{v4-arith:hecke:thm:weighted}
Fix $a\geq1$.
The same polynomial maps $c_{p^a}$ and $d_{p^a}$ extend boundedly
to every $\mathcal H_s$, with the following properties.
\begin{enumerate}
\item For $0\leq s<1/2$, $c_{p^a}$ is onto $\mathcal H_s$,
and $d_{p^a}$ is onto $\mathcal H_{s,\geq a}$.
Their common kernel has basis
\[
 g_\zeta=\sum_{r\geq0}(\zeta/p)^re_r,
 \qquad \zeta^{a+1}=1,\quad\zeta\ne1.
\]
Their polynomial inverses are not closable.
\item For $s=1/2$, both maps are injective.
Their ranges are dense, proper and nonclosed in $\mathcal H_s$
and $\mathcal H_{s,\geq a}$, respectively.
Their inverses on these actual ranges are closed and unbounded.
They are the graph closures of the polynomial inverses.
\item For $s>1/2$, $c_{p^a}$ is boundedly invertible on $\mathcal H_s$,
and $d_{p^a}$ is a bounded isomorphism onto $\mathcal H_{s,\geq a}$.
These inverses extend the polynomial inverses.
\end{enumerate}
For $a=0$, both maps are the identity for every $s$.
\end{theorem}
\begin{proof}
Set $b=p^{s+1/2}$, $q=p^{1/2-s}$ and $Q_a(z)=1+z+\cdots+z^a$.
The vectors $f_r=e_r/(u^{s+1/2}b^r)$ form an orthonormal basis.
Let $Vf_r=f_{r+1}$ and $L=V^*$.
The raw shifts become $U=bV$ and $pB=qL$, so
\begin{equation}
 \label{v4-arith:hecke:eq:weighted-shifts}
 c_{p^a}=Q_a(qL),\qquad d_{p^a}=b^aV^aQ_a(qL).
\end{equation}
This proves boundedness with bounds $\sum_{j=0}^aq^j$
and $b^a\sum_{j=0}^aq^j$.

The isometry $V^a$ identifies the space with its tail.
It therefore suffices to analyze $Q_a(qL)$.

For $q>1$, its kernel recurrence has distinct characteristic roots
$\zeta/q$, with $\zeta^{a+1}=1$ and $\zeta\ne1$.
All resulting geometric sequences are square-summable.
In raw coordinates, they are the stated $g_\zeta$, with
\[
 \|g_\zeta\|_s^2=\frac{u^{1+2s}}{1-p^{2s-1}}.
\]
The recurrence and Vandermonde argument prove that they form a basis.

Each $L-\zeta/q$ has bounded right inverse
$V(I-(\zeta/q)V)^{-1}$, by its Neumann series.
Factoring $Q_a(qL)$ and composing these right inverses in reverse
order gives a bounded right inverse $G$ for $c_{p^a}$.
Then
\[
 d_{p^a}(b^{-a}GL^a)=V^aL^a,
\]
the projection onto the tail.
This proves both surjectivity assertions.
Finite truncations of a nonzero kernel vector prove nonclosability
of the polynomial inverses, as above.

At $q=1$, kernel sequences are combinations of the periodic
sequences $\zeta^r$.
A nonzero periodic sequence is not square-summable, so the kernel vanishes.
The constant coefficient of $Q_a$ is one.
Its formal reciprocal acts on $D$, since $L$ is locally nilpotent there.
Thus $c_{p^a}(D)=D$, and the range is dense.

Fix one root $\zeta\ne1$ and, for $N>a$, set
\[
 v_N=N^{-1/2}\sum_{r=0}^{N-1}\zeta^rf_r.
\]
Its norm is one.
The first $N-a$ coefficients of $Q_a(L)v_N$ vanish.
At most $a$ coefficients remain, each bounded by $(a+1)/\sqrt N$.
Therefore
\[
 \|Q_a(L)v_N\|\leq\frac{\sqrt a(a+1)}{\sqrt N}\longrightarrow0.
\]

An injective bounded operator with closed range is bounded below,
by the bounded inverse theorem on its range.
The sequence excludes this possibility.
The dense range is consequently proper and nonclosed.
Equation~\eqref{v4-arith:hecke:eq:weighted-shifts} gives the same
conclusion inside the tail for $d_{p^a}$.

The inverse of an injective bounded operator on its range has closed
graph: interchange the coordinates of the original closed graph.
Here these inverses are unbounded by the preceding sequence.

For any $x\in\mathcal H_s$, choose $x_N\in D$ with $x_N\to x$.
Boundedness gives $c_{p^a}x_N\to c_{p^a}x$.
Hence $(c_{p^a}x_N,x_N)\to(c_{p^a}x,x)$ in the inverse graph.
The same argument applies to $d_{p^a}$.
This proves that the polynomial inverses have precisely these graph closures.

Finally, for $q<1$, the identity $(1-z)Q_a(z)=1-z^{a+1}$ gives
\[
 c_{p^a}^{-1}=(I-qL)(I-(qL)^{a+1})^{-1},\qquad
 \|c_{p^a}^{-1}\|\leq\frac{1+q}{1-q^{a+1}}.
\]
The inverse is a norm-convergent Neumann series.
On the tail, the inverse of $d_{p^a}$ is
$c_{p^a}^{-1}b^{-a}L^a$.
Both agree with the polynomial inverses on finite vectors.
\end{proof}

All these bounded maps retain the common finite domain $D$ as a core.
Their polynomial formulas do not change with $s$.
The completed kernels and inverses change because the allowed
coefficient decay changes.
In particular, excluding $g_\zeta$ from $\mathcal H_s$ at $s\geq1/2$
does not remove its kernel class in $\mathcal H_0$.

\section{Operations and states in the arithmetic comparison}
\label{v4-arith:hecke:sec:arithmetic}

The positive radial state belongs to its specified Hecke representation.
A phase average on oscillator states gives a different operator,
even when it uses the same measure.
To see this directly, take two species with polynomial space
$\mathbb C[a_k,b_k:k\geq1]$ and the Hermitian Hall norm.
The monomial $\prod_ka_k^{r_k}b_k^{s_k}$ has charge
$m=\sum_k(r_k-s_k)$.

Let $V_\theta$ multiply its charge space by $e^{im\theta}$.
These are unitary operators on the Hilbert completion.
The strong average
\[
 A_p=\int_0^\pi V_\theta\,d\nu_p(\theta)
\]
is a contraction. Its eigenvalue in charge $m$ is
$c_m=\int_0^\pi e^{im\theta}\,d\nu_p$.

Symmetry about $\pi/2$ gives
\[
 c_1=i\int_0^\pi\sin\theta\,d\nu_p\in i\mathbb R_{>0}.
\]
Thus this phase average cannot be unitarily intertwined with the
self-adjoint adjacency operator while retaining the specified operators.
The obstruction concerns this comparison, not every arithmetic realization.

The energy convention also matters.
For species labeled by primes, integer mode weight $k$ and
arithmetic energy $k\log p$ define different diagonal operators.
The integral monomial lattice survives both gradings,
although the real energy operator need not preserve it.

The full oscillator partition function is
$\prod_{k\geq1}\zeta(k\beta)$ for $\beta>1$.
Keeping only mode one instead gives $\zeta(\beta)$.
That subspace is not preserved by the index shifts or divisor substitutions.
The operator domains and direct occupation sums establishing these
statements are given in Section~\ref{v4-arith:ado:hilbert}.

For comparison, on $\ell^2(\mathbb N_{>0})$ let
$\mu_n\delta_N=\delta_{nN}$ and
$H_{\mathrm{BC}}\delta_N=\log N\,\delta_N$.
These commuting isometries satisfy
\[
 [H_{\mathrm{BC}},\mu_n]=\log n\,\mu_n
\]
on finite-support vectors.
They move $\delta_1$ to $\delta_n$.
The index substitutions instead fix their polynomial vacuum,
and their arithmetic energy satisfies $H\psi_n=n\psi_nH$.
A vacuum- and energy-preserving identification cannot equate these
two specified actions for $n>1$.

\begin{question}[Marked arithmetic realization]
\label{v4-arith:hecke:qn:arithmetic}
Construct an arithmetic trace carrier independently of these local models,
with its coefficient theory, integral form, selected trace class and functional.
Construct its operations and an intertwiner with a specified local
or restricted tensor-product representation.
The map must preserve the distinguished vector, the functional,
the determinant normalization and every chosen operation.
For a Hilbert comparison, specify the domains, adjoints and completed tensor product.
For a vertex comparison, retain translation and all singular products.
\end{question}

For example, the unramified comparison of
Conjecture~\ref{v4-arith:thm:beilinson-honest-reading} requires the
actual maps $\Phi,Q,Q_p$, normalized local vectors and compatible
linear functionals. An arbitrary family of Satake pairs does not
construct those arithmetic objects or maps.
Neither the root-state identity nor the Witt ghost formulas identify
a full Weil functional.

There is also an immediate operation obstruction to some proposed maps.
In characteristic zero, an injective intertwiner cannot identify
the divisor family with commuting Adams operations or a commuting
semigroup of isometries.
Equation~\eqref{v4-arith:hecke:eq:commutator} would map a nonzero vector to zero.
A quotient, a restricted domain or a different operation system remains
a separate possibility with its own kernel and comparison calculation.
